% GENERATED FILE. Edit canonical lessons, then run npm run generate:books.
% canonical-source-hash: fnv1a64:7fb94e00c73526e7
% canonical-lessons: IT-C534-lavastoviglie, IT-C534-sciacquare, IT-C534-spugna, IT-C534-lavandino, IT-C534-lucidare, IT-R534-ripasso-527-534

\chapter{Lavare i piatti}
\label{ch:it-lavare-i-piatti}
\hlchaptermodality{534}

\begin{chapteropening}
\textbf{By the end of this chapter:} \emph{I can talk about washing up and polishing.}

Complete the last lesson of chapter 534: I can talk about washing up and polishing.
\end{chapteropening}

\section[la lavastoviglie]{la lavastoviglie — dishwasher}

\label{lesson:IT-C534-lavastoviglie}



\begin{quote}
Before the new one: say the Italian for napkin, then the Italian for tray.
\end{quote}

\subsection*{You'll want to know: la lavastoviglie}
\textbf{la lavastoviglie} — \textquotedblleft{}dishwasher\textquotedblright{}.

\begin{quote}
Dopo cena carico la lavastoviglie.
\end{quote}

\begin{quote}
\emph{After dinner I load the dishwasher.}
\end{quote}

The same form in the plural.

\begin{cousinweb}
\emph{lavastoviglie} is built from \emph{lava} \textquotedblleft{}washes\textquotedblright{}, from Latin \emph{lavāre} (Lewis \& Short), and \emph{stoviglie} \textquotedblleft{}dishes\textquotedblright{}, whose origin is not settled.
\end{cousinweb}

\begin{grammarlens}[title={What you've built}]
One more word for telling, explaining and describing at B1.
\end{grammarlens}

\subsection*{Your turn}
\begin{itemize}
\raggedright
  \item \emph{Say it:} \emph{la lavastoviglie}
  \item \emph{Say it:} \emph{la lavastoviglie}, once more
  \item \emph{Say it:} a sentence of your own with \emph{la lavastoviglie}
  \item \emph{Recall:} say the Italian for napkin, then the Italian for tray, then say \emph{la lavastoviglie} again
\end{itemize}

\begin{tcolorbox}[breakable,colback=teal!4,colframe=teal!35!black,title={Before you move on}]
What does la lavastoviglie mean? (\textquotedblleft{}Dishwasher\textquotedblright{}.) Say it once more.
\end{tcolorbox}
\section[sciacquare]{sciacquare — to rinse}

\label{lesson:IT-C534-sciacquare}



\begin{quote}
Before the new one: say the Italian for tray, then the Italian for dishwasher.
\end{quote}

\subsection*{You'll want to know: sciacquare}
\textbf{sciacquare} — \textquotedblleft{}to rinse\textquotedblright{}.

\begin{quote}
Sciacquo i bicchieri sotto l'acqua.
\end{quote}

\begin{quote}
\emph{I rinse the glasses under the water.}
\end{quote}

\emph{Sciacquarsi la bocca} — to rinse one's mouth.

\begin{cousinweb}
\emph{sciacquare} is inherited from Late Latin \emph{exaquāre} (EtymDB), built from \emph{ex-} and \emph{aqua} \textquotedblleft{}water\textquotedblright{}.
\end{cousinweb}

\begin{grammarlens}[title={What you've built}]
One more word for telling, explaining and describing at B1.
\end{grammarlens}

\subsection*{Your turn}
\begin{itemize}
\raggedright
  \item \emph{Say it:} \emph{sciacquare}
  \item \emph{Say it:} \emph{sciacquare}, once more
  \item \emph{Say it:} a sentence of your own with \emph{sciacquare}
  \item \emph{Recall:} say the Italian for tray, then the Italian for dishwasher, then say \emph{sciacquare} again
\end{itemize}

\begin{tcolorbox}[breakable,colback=teal!4,colframe=teal!35!black,title={Before you move on}]
What does sciacquare mean? (\textquotedblleft{}To rinse\textquotedblright{}.) Say it once more.
\end{tcolorbox}
\section[la spugna]{la spugna — sponge}

\label{lesson:IT-C534-spugna}



\begin{quote}
Before the new one: say the Italian for dishwasher, then the Italian for to rinse.
\end{quote}

\subsection*{You'll want to know: la spugna}
\textbf{la spugna} — \textquotedblleft{}sponge\textquotedblright{}.

\begin{quote}
Lavo le pentole con la spugna.
\end{quote}

\begin{quote}
\emph{I wash the pans with the sponge.}
\end{quote}

\emph{Gettare la spugna} — to throw in the towel.

\begin{cousinweb}
\emph{spugna} is inherited from Latin \emph{spongia} \textquotedblleft{}a sponge\textquotedblright{} (EtymDB), from Greek \emph{spongiá}.
\end{cousinweb}

\begin{grammarlens}[title={What you've built}]
One more word for telling, explaining and describing at B1.
\end{grammarlens}

\subsection*{Your turn}
\begin{itemize}
\raggedright
  \item \emph{Say it:} \emph{la spugna}
  \item \emph{Say it:} \emph{la spugna}, once more
  \item \emph{Say it:} a sentence of your own with \emph{la spugna}
  \item \emph{Recall:} say the Italian for dishwasher, then the Italian for to rinse, then say \emph{la spugna} again
\end{itemize}

\begin{tcolorbox}[breakable,colback=teal!4,colframe=teal!35!black,title={Before you move on}]
What does la spugna mean? (\textquotedblleft{}Sponge\textquotedblright{}.) Say it once more.
\end{tcolorbox}
\section[il lavandino]{il lavandino — sink, washbasin}

\label{lesson:IT-C534-lavandino}



\begin{quote}
Before the new one: say the Italian for to rinse, then the Italian for sponge.
\end{quote}

\subsection*{You'll want to know: il lavandino}
\textbf{il lavandino} — \textquotedblleft{}sink, washbasin\textquotedblright{}.

\begin{quote}
Il lavandino della cucina è pieno di piatti.
\end{quote}

\begin{quote}
\emph{The kitchen sink is full of dishes.}
\end{quote}

Also \emph{il lavello} in the kitchen.

\begin{cousinweb}
\emph{lavandino} is built on \emph{lavanda} \textquotedblleft{}washing\textquotedblright{}, from Latin \emph{lavanda} \textquotedblleft{}things to be washed\textquotedblright{}, from \emph{lavāre} \textquotedblleft{}to wash\textquotedblright{} (Lewis \& Short).
\end{cousinweb}

\begin{grammarlens}[title={What you've built}]
One more word for telling, explaining and describing at B1.
\end{grammarlens}

\subsection*{Your turn}
\begin{itemize}
\raggedright
  \item \emph{Say it:} \emph{il lavandino}
  \item \emph{Say it:} \emph{il lavandino}, once more
  \item \emph{Say it:} a sentence of your own with \emph{il lavandino}
  \item \emph{Recall:} say the Italian for to rinse, then the Italian for sponge, then say \emph{il lavandino} again
\end{itemize}

\begin{tcolorbox}[breakable,colback=teal!4,colframe=teal!35!black,title={Before you move on}]
What does il lavandino mean? (\textquotedblleft{}Sink, washbasin\textquotedblright{}.) Say it once more.
\end{tcolorbox}
\section[lucidare]{lucidare — to polish, to shine}

\label{lesson:IT-C534-lucidare}



\begin{quote}
Before the new one: say the Italian for sponge, then the Italian for sink.
\end{quote}

\subsection*{You'll want to know: lucidare}
\textbf{lucidare} — \textquotedblleft{}to polish, to shine\textquotedblright{}.

\begin{quote}
Ho lucidato le scarpe per il matrimonio.
\end{quote}

\begin{quote}
\emph{I polished my shoes for the wedding.}
\end{quote}

Built on \emph{lucido} \textquotedblleft{}shiny\textquotedblright{}.

\begin{cousinweb}
\emph{lucidare} is built on \emph{lucido} \textquotedblleft{}shiny\textquotedblright{}, a learned borrowing from Latin \emph{lūcidus} \textquotedblleft{}bright\textquotedblright{} (Lewis \& Short), from \emph{lūx} \textquotedblleft{}light\textquotedblright{}.
\end{cousinweb}

\begin{grammarlens}[title={What you've built}]
One more word for telling, explaining and describing at B1.
\end{grammarlens}

\subsection*{Your turn}
\begin{itemize}
\raggedright
  \item \emph{Say it:} \emph{lucidare}
  \item \emph{Say it:} \emph{lucidare}, once more
  \item \emph{Say it:} a sentence of your own with \emph{lucidare}
  \item \emph{Recall:} say the Italian for sponge, then the Italian for sink, then say \emph{lucidare} again
\end{itemize}

\begin{tcolorbox}[breakable,colback=teal!4,colframe=teal!35!black,title={Before you move on}]
What does lucidare mean? (\textquotedblleft{}To polish, to shine\textquotedblright{}.) Say it once more.
\end{tcolorbox}
\section[(dialogue)]{Ripasso dei capitoli 527–534}

\label{lesson:IT-R534-ripasso-527-534}



\begin{quote}
Say the Italian for biography and science fiction.
\end{quote}

\subsection*{Your turn}
Say these aloud:

\begin{itemize}
\raggedright
  \item biography, science fiction, essay, comic, poet
  \item series, instalment, episode, spoiler, remote control
  \item to dub, dubbing, subtitles, soundtrack, trailer
  \item to subscribe, streaming, video game, to post, follower
  \item laundry, to hang, clothes horse, clothes peg, iron
  \item to sweep, to dust, rag, detergent, vacuum cleaner
  \item to lay, to clear, tablecloth, napkin, tray
  \item dishwasher, to rinse, sponge, sink, to polish
\end{itemize}

\begin{tcolorbox}[breakable,colback=teal!4,colframe=teal!35!black,title={Before you move on}]
Biography? (\textbf{la biografia}.) Science fiction? (\textbf{la fantascienza}.) Essay? (\textbf{il saggio}.) Comic? (\textbf{il fumetto}.) Poet? (\textbf{il poeta}.) Series? (\textbf{la serie}.) Instalment? (\textbf{la puntata}.) Episode? (\textbf{l'episodio}.) Spoiler? (\textbf{lo spoiler}.) Remote control? (\textbf{il telecomando}.) To dub? (\textbf{doppiare}.) Dubbing? (\textbf{il doppiaggio}.) Subtitles? (\textbf{i sottotitoli}.) Soundtrack? (\textbf{la colonna sonora}.) Trailer? (\textbf{il trailer}.) To subscribe? (\textbf{abbonarsi}.) Streaming? (\textbf{lo streaming}.) Video game? (\textbf{il videogioco}.) To post? (\textbf{postare}.) Follower? (\textbf{il follower}.) Laundry? (\textbf{il bucato}.) To hang? (\textbf{appendere}.) Clothes horse? (\textbf{lo stendino}.) Clothes peg? (\textbf{la molletta}.) Iron? (\textbf{il ferro da stiro}.) To sweep? (\textbf{spazzare}.) To dust? (\textbf{spolverare}.) Rag? (\textbf{lo straccio}.) Detergent? (\textbf{il detersivo}.) Vacuum cleaner? (\textbf{l'aspirapolvere}.) To lay? (\textbf{apparecchiare}.) To clear? (\textbf{sparecchiare}.) Tablecloth? (\textbf{la tovaglia}.) Napkin? (\textbf{il tovagliolo}.) Tray? (\textbf{il vassoio}.) Dishwasher? (\textbf{la lavastoviglie}.) To rinse? (\textbf{sciacquare}.) Sponge? (\textbf{la spugna}.) Sink? (\textbf{il lavandino}.) To polish? (\textbf{lucidare}.)
\end{tcolorbox}
